\documentclass[10pt,a4paper]{amsart}
\usepackage[margin=19mm]{geometry}
\usepackage{amsmath,amssymb,mathtools}
\usepackage[scr=rsfs,cal=euler]{mathalfa}
\usepackage{xfrac}
\usepackage[unicode,hidelinks,bookmarks=false]{hyperref}
\newcommand{\ie}{i.e.\ }
\newcommand{\cf}{cf.\ }
\newcommand{\resp}{respectively\ }
\newcommand{\Subsection}{\subsection}
\providecommand{\nobreakdash}{\nobreak\hbox{-}}
\setlength{\emergencystretch}{2em}
\multlinegap0pt
\usepackage{bm}
\usepackage{bbm}
\usepackage{dsfont}
\usepackage{xspace}
\usepackage{cancel}
\usepackage{mathtools}
\usepackage{cleveref}
\usepackage{amsthm}
\makeatletter
\@ifundefined{theorem}{\newtheorem{theorem}{Theorem}}{}
\@ifundefined{lemma}{\newtheorem{lemma}[theorem]{Lemma}}{}
\@ifundefined{proposition}{\newtheorem{proposition}[theorem]{Proposition}}{}
\makeatother
\makeatletter
\DeclareMathOperator{\Adj}{Adj}
\DeclareMathOperator{\adj}{adj}
\newcommand{\ci}{\mathcal{I}}
\expandafter\ifx\csname example\endcsname\relax
\newenvironment{example}
  {\pushQED{\qed}\renewcommand{\qedsymbol}{$\triangle$}\examplex}
  {\popQED\endexamplex}
\fi
\newcommand{\zz}{\mathbb{Z}}
\makeatother
\title[Parke-Taylor varieties]{Parke-Taylor varieties}
\author{Benjamin Hollering, Dmitrii Pavlov}
\date{Source: arXiv:2509.09323v1 (CC BY 4.0)}
\begin{document}
\maketitle

\noindent\emph{Source and licence.} Parke-Taylor varieties. Version arXiv:2509.09323v1 (CC BY 4.0), distributed under the Creative Commons Attribution 4.0 International licence, as recorded in the arXiv licence metadata for this exact version and shown on the abstract page https://arxiv.org/abs/2509.09323v1. Full rights notice: licenses/arxiv-2509.09323-CC-BY-4.0.txt.\par\smallskip

\noindent\textit{Editorial scope.} Counted unit: the Lemma whose source label is \texttt{lemma:lift-binomials} in the article (source0.tex lines 807--824, source label \texttt{lemma:lift-binomials}), delivered together with the article's own complete original proof of it (source0.tex lines 825--846, one \texttt{proof} environment of 22 lines). No mathematics is written, repaired or re-derived: statement and proof are the article's own text line for line. Required context retained uncounted: prop:toric-gens-via-adjs (lines 689--694). Every definition, display and label of those lines is retained unchanged. Omitted from this package and not redistributed: 1--688, 695--806, 847--1108. Nothing inside a retained excerpt is omitted or altered: every retained body is delivered exactly as the article's lines are written, so no omission receipt of any kind accompanies this package. Citations: the retained excerpts carry no citation command at all, so no bibliography entry of the article is transcribed and no reference is redistributed. Reference adaptation: none was needed and none was made; every reference inside the delivered unit resolves inside it, so no unresolved reference or citation is produced. Printed slips kept literal: no printed word, symbol, hypothesis or number of the counted statement or of its proof was altered. Layout and class: the article's own class is \texttt{amsart} with packages including inputenc, babel, ifthen, microtype, xcolor, comment, blkarray, fullpage, mathtools, xfrac, shuffle, centernot, enumitem, mathdots; the wrapper is the lane's pinned \texttt{amsart} editorial wrapper at 10pt on A4, which restates the theorem-like environments the retained text needs and adds the article's own macro definitions that the retained text needs (\texttt{\textbackslash Adj}, \texttt{\textbackslash adj}, \texttt{\textbackslash ci}, \texttt{\textbackslash zz}), with extra wrapper packages (bm, bbm, dsfont, xspace, cancel, mathtools, cleveref). The delivered numbering is the unit's own. Assets: no retained span contains a figure environment, a graphics-inclusion command, a PDF-page inclusion, a table or a TikZ picture; no image file is copied into this package, the package declares no asset dependency and no image file is redistributed with it. Licence: version arXiv:2509.09323v1 of this article is distributed under the Creative Commons Attribution 4.0 International licence, as recorded in the arXiv licence metadata for this exact version and shown on the abstract page https://arxiv.org/abs/2509.09323v1. A full rights notice is delivered at licenses/arxiv-2509.09323-CC-BY-4.0.txt. Characters: every retained excerpt is pure ASCII, so the delivered engine, XeLaTeX with its default Latin Modern text font, reports no missing character.
\begin{proposition}
\label{prop:toric-gens-via-adjs}
A binomial equation of the form 
$$\prod\limits_{\sigma\in A} z_\sigma - \prod\limits_{\tau\in B} z_\tau = 0$$
holds on $T_n$ if and only if the set of adjacencies of all permutations in $A$ (counted with multiplicity) is the same as the set of adjacencies of all permutations in $B$. Moreover, the set of all such binomials above generates $\ci(T_n)$. 
\end{proposition}

\begin{lemma}
\label{lemma:lift-binomials}
Suppose that $A, B \subset \Sigma_n$ such that $\prod\limits_{\sigma\in A} z_\sigma - \prod\limits_{\tau\in B} z_\tau \in \ci(T_n)$ and $|A| = |B| = m$. If there exists an index $i \in \{2,\ldots, n\}$ such that $\{\sigma_i \sigma_{i+1} ~:~ \sigma \in A \} = \{\tau_i \tau_{i+1} ~:~ \tau \in B \}$ as multisets, then for any $k \in \zz_{>0}$, it holds that
\begin{equation}
\label{eqn:lift-binomial}
\begin{bmatrix}
1 & 2 & \cdots & \sigma_i^{(1)} & \delta & \sigma_{i + 1}^{(1)} & \cdots & \sigma_n^{(1)} \\
1 & 2 & \cdots & \vdots         & \delta & \vdots               & \cdots & \vdots  \\
1 & 2 & \cdots & \sigma_i^{(m)} & \delta & \sigma_{i + 1}^{(m)} & \cdots & \sigma_n^{(m)} \\
\end{bmatrix} - 
\begin{bmatrix}
1 & 2 & \cdots & \tau_i^{(1)} & \delta & \tau_{i + 1}^{(1)} & \cdots & \tau_n^{(1)} \\
1 & 2 & \cdots & \vdots       & \delta & \vdots             & \cdots & \vdots  \\
1 & 2 & \cdots & \tau_i^{(m)} & \delta & \tau_{i + 1}^{(m)} & \cdots & \tau_n^{(m)} \\
\end{bmatrix} \in \ci(T_{n + k}), 
\end{equation}
where $\delta = \delta_1 \delta_2 \cdots \delta_k$ is a permutation on the letters $n+1, \ldots, n+k$.
\end{lemma}

\begin{proof}
The equality $\{ \sigma_i \sigma_{i+1} ~:~ \sigma \in A \} = \{\tau_i \tau_{i+1}~:~ \tau \in B \}$ of the multiset of adjacencies at position $i$ implies 
that the multisets
\[
\{ \sigma_i \delta_1,~ \delta_k \sigma_{i+1} ~~:~ \sigma \in A \} = \{ \tau_i \delta_1,~ \delta_k \tau_{i+1} ~:~ \tau \in B\}
\]
are also equal. For any permutation $\pi \in \Sigma_n$, let $\Adj(\pi)$ denote the set of adjacencies of $\pi$ taken cyclically and $\adj(\pi)$ denote the standard set of adjacencies (meaning taken acyclically). We have
\begin{equation}
\label{eqn:lift-perm-adj-decomp}
\Adj(12 \sigma_3 \cdots \sigma_i \delta \sigma_{i+1} \ldots \sigma_n) = \adj(12\ldots \sigma_i) \cup\{\sigma_i \delta_1, \delta_k \sigma_{i+1
} \} \cup \adj(\sigma_{i+1} \ldots \sigma_n 1).   
\end{equation}


Since $\prod_{\sigma\in A} z_\sigma - \prod_{\tau\in B} z_\tau \in \ci(T_n)$, by \cref{prop:toric-gens-via-adjs}, it holds that $\cup_{\sigma \in A} \Adj(\sigma) = \cup_{\tau \in B} \Adj(\tau)$ and thus by assumption it follows that $\cup_{\sigma \in A} \Adj(\sigma) \setminus \{\sigma_i \sigma_{i+1}\} = \cup_{\tau \in B} \Adj(\tau) \setminus \{\tau_i \tau_{i+1}\}$. Combining this equality together with the fact that
\begin{align}
\label{eqn:perm-union-decomp}
&\cup_{\sigma \in A} \Adj(\sigma) \setminus \{\sigma_i \sigma_{i+1} \} = \cup_{\sigma \in A} \adj(\sigma_1 \ldots \sigma_i) \cup \adj(\sigma_{i+1} \ldots \sigma_n 1), \\
&\cup_{\tau \in B} \Adj(\tau) \setminus \{\tau_i \tau_{i+1}\} = \cup_{\tau\in A} \adj(\tau_1 \ldots \tau_i) \cup \adj(\tau_{i+1} \ldots \tau_n 1)
\end{align}
yields that the multisets of adjacencies of the permutations in the two monomials in \cref{eqn:lift-binomial} are the same and thus by \cref{prop:toric-gens-via-adjs} the claim holds. 
\end{proof}

\end{document}
